% ================================================================
%  chapters/parts/ch01-b.tex  —  第 1 章 分片 b
%  §1.2.2  孔容和孔尺寸分布
%  源：S1（上册）§1.2.2，书内 p.43—60 / PDF p.62—79
%  接 ch01-a.tex（比表面积）之后；颗粒度、密度、机械强度见分片 c
% ================================================================

\subsection{孔容和孔尺寸分布}

\begin{probout}
\textbf{本节回答}：一条已经测好的吸附等温线，如何变成孔容与孔尺寸分布，这些数值在什么
条件下可信，以及为什么不同方法报出的「平均孔径」不能直接互相比较。比表面积的求算流程
见 \chapref{ch:01} 前一节。

\textbf{输出}：比孔容（$\unit{cm^3.g^{-1}}$ 或 $\unit{mL.g^{-1}}$）与孔尺寸分布曲线
——通常画成 $\mathrm{d}V_p/\mathrm{d}\lg D_p$ 对 $\lg D_p$，以及由分布派生的总孔容、
最可几孔宽和外表面积。

\textbf{本节核心判断}：孔尺寸分布\textbf{不是被「测」出来的，而是用一个模型从等温线
反演出来的}。原始数据只有一条曲线；要把它变成分布，必须补上三重假定：（i）孔的形状是
圆柱形、平板狭缝还是球形；（ii）吸附层厚度 $t$ 与相对压力、与孔壁曲率的关系；
（iii）孔之间互不连通、互不干扰。这三条都写在模型里，不在数据里。于是\textbf{换模型就
换分布}：同一组数据用 BJH 与 NLDFT 反演，峰的位置、个数和高度都可以不同
（图~\ref{fig:psd-model}）。因此报告孔尺寸分布时\textbf{必须同时给出实验方法、吸附质和
模型（含 kernel）名称}三项，缺任何一项，数值都无法复核，也无法与文献比较。
\end{probout}

\subsubsection{孔、孔宽与等效孔模型}

按与外界是否连通，孔分为\textbf{开孔}（open pore，与外界连通的孔）、\textbf{盲孔}
（blind pore，只有一个外界连通开口的开孔）与\textbf{闭孔}（closed pore，不与外界连通的
孔）。只有反应物、产物分子能进出的开孔对催化有意义；闭孔不参与吸附与扩散，气体吸附法
测不到它，需用小角 X 射线散射或小角中子散射，本节不涉及。

孔尺寸用\textbf{孔宽}（pore width）标示：圆柱形孔取孔直径而非半径，平板孔取板壁间距，
二者统称孔宽。作这一约定是为了让两种等效孔模型的数值对吸附质的容纳能力具有可比性；本书
一律用「孔尺寸」而非「孔径」，因为实际的孔极少是圆的。

真实催化剂内部的孔极不规则，无法用纯几何特征值描述。间接测定法所做的，是把真实孔系统
\textbf{等效为一组尺寸不同、形状规整、互不连通的模型孔}，使流体在这组模型孔上的渗透与
滞留行为与真实样品相同。由此得到的分布是\textbf{等效模型孔的理想分布，只是名义值}，
具有平均或等效的意义，不再反映真实孔的实际几何；它\textbf{完全依赖于测定方法（流体
分子的性质）和分析理论（所采用的孔模型）}。虽然失去了确定的几何意义，这一结果对理解
催化剂的吸附能力与催化活性恰恰十分有用。还须注意，现有孔分析理论都假定\textbf{各孔互不
连通、相互独立}，而真实催化剂是多孔网络，「墨水瓶」孔、并行喉管、连通腔体都超出模型的
表达能力；孔之间的连通性目前只能用回滞环扫描等手段定性了解
\src{S1 §1.2.2，书内 p.43--45 / PDF p.62--64}。

\subsubsection{毛细凝聚与 Kelvin 方程}

\begin{principle}
\textbf{物理图像}：中孔内先由吸附质以液膜形式覆盖孔壁，液膜厚度 $t$ 随 $\ppzero$ 增大而
增厚；当液膜围成的「孔心」达到弯曲液面的平衡条件时，孔心发生\textbf{毛细凝聚}
（capillary condensation），孔被凝聚液充满。描述这一平衡的是 Kelvin 方程：

\begin{equation}\label{eq:kelvin}
\ln\frac{p}{p_0}=-\frac{2\gamma V_{\mathrm{L}}\cos\theta}{r_kRT}
\end{equation}

式中 $\gamma$ 为液体的表面张力，$V_{\mathrm{L}}$ 为液体的摩尔体积，$\theta$ 为液体与孔壁
的接触角（氮在 $77\,\unit{K}$ 下润湿孔壁，$\cos\theta\approx1$），$R$ 为摩尔气体常量，
$T$ 为热力学温度，$r_k$ 为\textbf{弯月面的曲率半径}，称 Kelvin 半径。

量纲检查：$\gamma V_{\mathrm{L}}$ 为 $\unit{N.m^{-1}}\times\unit{m^3.mol^{-1}}
=\unit{N.m^2.mol^{-1}}$，除以 $r_k$（$\unit{m}$）得 $\unit{N.m.mol^{-1}}$，与
$RT$（$\unit{J.mol^{-1}}=\unit{N.m.mol^{-1}}$）之比无量纲，与等式左边的对数一致。

\textbf{$r_k$ 不是孔半径}：它只是发生凝聚的孔心半径，而孔壁上已存在厚度为 $t$ 的吸附层，
故真实孔半径 $r_p=r_k+t$。这一步称\textbf{曲率修正}；不做修正而直接以 $r_k$ 当孔半径，
会把孔算小，并把表面积算大。
\end{principle}

\begin{keycond}
\textbf{（1）弯月面形状决定系数。}式~\eqref{eq:kelvin} 中的系数 $2$ 对应\textbf{半球形}
弯月面；对圆柱形弯月面，平均曲率减半，系数应改为 $1$。脱附时凝聚液自孔口以半球形弯月面
蒸发，吸附时孔内凹液面为圆柱面，二者给出不同的 $r_k$--$\ppzero$ 关系——这正是回滞环的
几何来源。使用前必须先确认所用公式对应吸附支还是脱附支。
\textbf{（2）只在中孔内成立。}微孔的孔径与吸附质分子尺寸同量级，孔壁势场叠加，流体不再
具有体相性质，连续介质图像失效。
\textbf{（3）数值常数与吸附质、温度绑定。}以 $77\,\unit{K}$ 氮吸附为例，取
$\gamma=\qty{8.85}{mN.m^{-1}}$、$V_{\mathrm{L}}=\qty{34.7}{cm^3.mol^{-1}}$、
$\cos\theta=1$，则 $2\gamma V_{\mathrm{L}}/(RT)=\qty{0.955}{nm}$，即
$r_k/\unit{nm}=0.955\big/\ln(p_0/p)$；代入 $\ppzero=0.5$ 得 $r_k\approx\qty{1.38}{nm}$，
同点 $t\approx\qty{0.68}{nm}$，故 $r_p\approx\qty{2.06}{nm}$、$D_p\approx\qty{4.1}{nm}$。
该算例用于校核常数取值，不是实验值。\tocheck{未对照原页面图像：原书在 Kelvin 公式处给出的
常数取值与量纲约定（$r_k$ 以 nm 或 Å、是否已含 $\cos\theta$）OCR 未保留，上列常数系按
标准物理量自行计算。}
\end{keycond}

\subsubsection{BJH 法：中孔分布的记账框架}

BJH 法（Barrett--Joyner--Halenda 法）采用\textbf{圆柱形盲孔等效模型}，假定孔之间互不
连通、各孔在吸附或脱附过程中彼此独立，并假定中孔内液膜厚度 $t$ \textbf{只与相对压力
有关}，与吸附质性质和孔半径无关。其核心是一步守恒记账：相对压力 $x_k=(p/p_0)_k$ 下的
总吸附量等于\textbf{已凝聚孔内的液量}与\textbf{未凝聚孔壁上的液膜量}之和。

\begin{equation}\label{eq:bjh-balance}
V(x_k)=\underbrace{\sum_{i\le k}\pi r_i^{2}\ell_i}
_{\text{半径 }r_i\le r_k\text{ 的孔：已充满，计全孔体积}}
\;+\;\underbrace{\sum_{j>k}\pi\!\left[r_j^{2}-\left(r_j-t_k\right)^{2}\right]\ell_j}
_{\text{尚未凝聚的孔：只计孔壁液膜体积}}
\end{equation}

式中 $r_i$、$\ell_i$ 为第 $i$ 组孔的半径与该组孔总长度，$r_k$ 为 $x_k$ 对应的 Kelvin
半径，$t_k=t(x_k)$ 为对应液膜厚度，且孔半径与膜厚满足 $r_i=r_k(x_i)+t(x_i)$——孔壁上先有
厚度 $t$ 的吸附层，余下的孔心半径才是 Kelvin 半径（见本节前文曲率修正）。
\textbf{两项记的是两件不同的事，不能互相套用公式}：
\textbf{（1）已凝聚的孔}（$r_i\le r_k$）内是\textbf{充满}的液体，单位孔长的液量就是全孔
截面积 $\pi r_i^{2}$，即整个体积 $\Delta V_i=\pi r_i^{2}\ell_i$；
\textbf{（2）尚未凝聚的孔}（$r_j>r_k$）只有孔壁上的液膜，单位孔长液量为全截面积减去
未填的孔心 $\pi[r_j^{2}-(r_j-t_k)^{2}]$，即等价的环形面积。若把第 (1) 项误写成
$\pi[r_i^{2}-(r_i-t_k)^{2}]$，等于把已充满的孔当成只剩一圈液膜的孔：取 $r_i=\qty{5}{nm}$、
$t_k=\qty{1}{nm}$，正确的单位孔长液量是 $25\pi\,\unit{nm^2}$，误式只给 $9\pi\,\unit{nm^2}$
（少 $64\%$）；更荒谬的是令 $t_k\to0$ 时误式给出\textbf{零体积}，而充满的孔
体积与膜厚无关。

\textbf{液膜项的薄膜近似}：环面积可展开为
$r_j^{2}-(r_j-t_k)^{2}=2r_jt_k-t_k^{2}$，当 $t_k\ll r_j$ 时略去 $t_k^{2}$ 得
$2r_jt_k$（$t_k/r_j=0.04$ 时高估 $2\%$）。本书保留精确的环形写法，只有在孔半径远大于
膜厚、且明确注明薄膜近似时才用 $2\pi r_j\ell_jt_k$。

式~\eqref{eq:bjh-balance} 各项都是体积，故可直接相加；
相邻两个相对压力下的差值 $V(x_{k+1})-V(x_k)$ 给出 $r_k$ 至 $r_{k+1}$ 区间的孔容增量，
累加得积分分布，求导得微分分布 $\mathrm{d}V_p/\mathrm{d}D_p$，其最高峰对应
\textbf{最可几孔宽}。

\textbf{递推的物理含义}：相对压力升高时，孔壁液膜继续增厚，毛细凝聚同时向更大的孔推进，
两个过程都在消耗吸附量，必须逐级扣除，不能把吸附量增量全部记到新凝聚的那组孔上——这
正是 BJH 手工推算冗长繁琐的原因。石油化工行业标准 SH/T 0572—1993《孔尺寸分布计算法
（氮脱附等温线计算法）》给出这套递推的具体算法，商品吸附仪均内置相应程序；孔宽与相对
压力的对应关系即

\begin{equation}\label{eq:bjh-dp}
D_p=2\left(r_k+t\right),\qquad \ppzero=0.95\;\Longleftrightarrow\;D_p\approx\qty{50}{nm}
\end{equation}

\tocheck{未对照原页面图像：原书式(1-36)的符号体系（$V(x_k)$、$r_i$、$t_k$ 的下标约定与
求和上下限）OCR 完全失真，式~\eqref{eq:bjh-balance} 按 BJH 原始记账方式重写。其中
\textbf{第 (1) 项「已充满孔计全孔体积 $\pi r_i^{2}\ell_i$」不是从 OCR 判读出来的，而是由
该式自身的物理含义推定的}——已凝聚的孔内是充满的液体，写成环形面积
$\pi[r_i^{2}-(r_i-t_k)^{2}]$ 会让「充满的孔」在 $t_k\to0$ 时体积归零，与该式的定义自相
矛盾（本节已就该式加数值自检）。原书该项的确切写法仍需查 PDF p.65 原页面核对；
SH/T 0572 的标准年份与完整题名亦未在图像上确认。}

\begin{keycond}
\textbf{BJH 的三个硬边界。}（1）\textbf{不适用于微孔}：$\ppzero<0.2$ 的吸附量变化由微孔
填充贡献，Kelvin 方程在此无物理意义，强行计算会得到虚假峰。（2）\textbf{孔宽有上下限}：
由式~\eqref{eq:bjh-dp}，$77\,\unit{K}$ 氮吸附下 BJH 计算中孔孔宽的\textbf{上限约
$\qty{50}{nm}$}，用脱附支计算时可靠\textbf{下限为 $\qtyrange{5}{6}{nm}$}；两端之外应改用
压汞法（大孔端）或微孔方法（小孔端）。（3）\textbf{孔宽小于临界宽度 $D_{\mathrm{ch}}$ 时
没有回滞环}：此时等温线为完全可逆的 IV(b) 型，但仍可由毛细凝聚造成台阶。孔宽约
$\qty{4}{nm}$ 的 MCM-41 在 $77\,\unit{K}$ 氮吸附下即为无回滞环的 IV(b) 型，台阶跃升位于
$\ppzero=0.41$--$0.46$，由柱状孔几何算得的平均孔宽 $\qty{3.2}{nm}$ 仍真实反映了该材料；
换用 $77\,\unit{K}$ 氧吸附后等温线变为 IV(a) 型，回滞环出现在 $\ppzero=0.34$--$0.44$，
估算孔宽 $\qtyrange{3.0}{3.2}{nm}$。\textbf{无回滞环不等于无中孔}
\src{S1 §1.2.2 1)，书内 p.45--48 / PDF p.64--67}。
\end{keycond}

BJH 假定孔内液膜的化学势等于液体本体的化学势，但这样的液膜在热力学上是不稳定的。
Broekhoff 与 de Boer 考虑了两者化学势的差别，指出液膜厚度不仅与相对压力有关，还与液膜
曲率（即孔半径）有关，据此给出满足热力学平衡条件的 \textbf{BdB 法}；当回滞环明显、孔壁
曲率影响不可忽略时，BdB 法比 BJH 法更自洽，代价是计算量更大。

\subsubsection{$t$ 图法与微孔分析：MP 法}

\begin{principle}
\textbf{前提}：非孔性固体上，当表面吸附多于单层后，氮的吸附层数 $n$（或厚度 $t$）
\textbf{只与相对压力有关，而与吸附剂种类无关}，因此对一种吸附质存在一条\textbf{通用标准
吸附等温线}（$t$ 曲线）。这是全部 $t$ 类方法的共同基石，它的失效方式也决定了这些方法的
误差来源。吸附层厚度与相对压力的经验关系（Halsey 方程）为

\begin{equation}\label{eq:halsey}
t=t_{\mathrm{m}}\left[\frac{-5}{\ln\left(\ppzero\right)}\right]^{1/3}
\end{equation}

$t_{\mathrm{m}}$ 为氮的单分子层平均厚度，取 $\qty{0.354}{nm}$；式~\eqref{eq:halsey}
在 $\ppzero\to1$ 时发散，只能用于毛细凝聚显著之前的区间。
\tocheck{未对照原页面图像：原书式(1-37) 的指数与常数（$-5$、$1/3$）在 OCR 中失真，
此处按 Halsey 方程标准形式重建；$t_{\mathrm{m}}=\qty{0.354}{nm}$ 取自 OCR 正文。}吸附量与厚度成正比：

\begin{equation}\label{eq:vt}
V=\frac{V_{\mathrm{m}}}{t_{\mathrm{m}}}\,t
\end{equation}

即 $V$--$t$ 图在非孔性吸附剂上应是一条\textbf{过原点的直线}，斜率等于
$V_{\mathrm{m}}/t_{\mathrm{m}}$，说明吸附量增加时留存表面积不变。\textbf{若吸附量以
凝聚态液氮体积为单位，则该斜率直接给出比表面积}：斜率单位为
$\unit{cm^3.g^{-1}}/\unit{nm}=10^{7}\,\unit{cm^2.g^{-1}}$，而 $\qty{1}{cm^3}$ 液氮铺成
厚 $\qty{0.354}{nm}$ 的膜时覆盖 $1/0.354\times10^{7}\approx2.8\times10^{7}\,
\unit{cm^2}$，两者一致，故有换算
$S/(\unit{m^2.g^{-1}})\approx10^{3}\times\dfrac{V_{\text{液}}/(\unit{cm^3.g^{-1}})}
{t/(\unit{nm})}$；用标准状况体积表示时等价于 $S\approx1.54\,V_{\mathrm{STP}}/t$。
\tocheck{未对照原页面图像：原书在 $V$--$t$ 斜率处只写「此斜率即为样品的表面积」，未标明
体积单位与 $10^{3}$ 这一换算因子；上列因子按液氮密度自行推导。}
\end{principle}

\textbf{多孔吸附剂的 $V$--$t$ 曲线形状与判读。}与过原点直线相比共有四种典型偏离方式：
\textbf{微孔}存在时，很低的相对压力下即发生微孔填充，被充满的微孔不能再发展多层吸附、
留存表面积下降，曲线\textbf{向下（上凸）偏离}；\textbf{中孔}存在时，中等相对压力下发生
毛细凝聚，吸附量急剧增加而标准 $t$ 值增加很小，曲线\textbf{向上偏离}；\textbf{中孔与微孔
共存}时曲线呈 \textbf{S 形}；毛细凝聚完成后吸附只在外表面继续，曲线\textbf{再次成为
直线}，其\textbf{斜率反映外表面积}、\textbf{截距等于中孔孔容}
\src{S1 §1.2.2 3)a，书内 p.52--53 / PDF p.71--72}。

\begin{keycond}
\textbf{$t$ 方法的两条实用条件。}（1）$V$--$t$ 曲线\textbf{并不总能过原点}，这是吸附质
与吸附剂相互作用强度不同造成的；按化学性质或按相互作用强度分级建立多套标准等温线，正是
对 $t$ 方法的改进方向。（2）用低相对压力段（$\ppzero<0.4$ 附近）的直线斜率求总表面积、
用高相对压力段的直线截距求中孔孔容，是\textbf{两个彼此独立的步骤}，不能用同一段直线同时
得出；进入回滞区后 $t$ 值物理意义消失，两段直线的判读随之失效。
\end{keycond}

\textbf{MP 法（微孔法）}是 $t$ 方法在微孔上的扩展。微孔吸附剂的 $V$--$t$ 曲线必然向下
弯曲偏离直线（反之不成立），下弯段本身即携带微孔孔尺寸信息。取平行板狭缝模型，微孔容
$V$、孔壁面积 $S$ 与厚度 $t$ 满足 $V=S\,t$；在弯曲段某点 $t_i$ 处的\textbf{切线斜率即
该厚度下留存的表面积} $S_i$，厚度增至 $t_{i+1}$ 时更大的微孔被充满，而被充满那组微孔的
板间距正对应 $t_{i+1}$。于是第 $i$ 组微孔的孔容为

\begin{equation}\label{eq:mp-dv}
\Delta V_i=\Delta S_i\,\bar{t}_i,\qquad
\Delta S_i=S_{i-1}-S_i,\qquad
\bar{t}_i=\tfrac12\left(t_i+t_{i+1}\right)
\end{equation}

由 $\Delta V_i$ 对该组孔宽（平板模型取 $D_{p,i}\approx2t_i$）作图即得微孔孔尺寸分布。
量纲检查：$\Delta S_i$ 为 $\unit{m^2.g^{-1}}$、$\bar{t}_i$ 为 $\unit{nm}$，乘积为
$\unit{nm.m^2.g^{-1}}=10^{-7}\,\unit{cm^3.g^{-1}}$，与孔容单位一致。
\tocheck{未对照原页面图像：原书式(1-42)--(1-44) 的下标约定（$\bar t_i$ 取相邻均值还是端点
值）与「板间距 $=t_{i+1}$ 还是 $=2t_i$」在 OCR 中无法判读，此处按 MP 法通用约定重写。}
两个使用要点：其一，低 $t$ 段过原点直线的斜率给出\textbf{总表面积}，大 $t$ 段直线外推的
\textbf{截距给出微孔总孔容}、斜率给出\textbf{外表面积}；该分析\textbf{隐含平板狭缝孔
假设，不适用于圆柱形微孔}，且微孔材料的总表面积本身物理意义有限。其二，MP 法同样按 BET
公式的 $C$ 值划分相互作用强弱并选择相应 $t$ 曲线（同一相对压力下相互作用强的体系液膜
稍厚），也可用与样品化学性质相同的非孔性吸附剂作基准物建立专属标准等温线
\src{S1 §1.2.2 3)c，书内 p.54--56 / PDF p.73--75}。

\subsubsection{标准吸附等温线对比法：$t$、$\alpha_s$ 与 $n$ 曲线}

$t$ 方法、$\alpha_s$ 方法与 $n$ 方法\textbf{都不是从理论模型推导出来的}，而是从大量
实验数据中归纳出的经验方法，同属「标准吸附等温线对比法」。它们的共同操作是：把多孔样品
的吸附量 $V$ 对一个描述吸附层状态的归一化量作图，由直线的斜率与弯曲位置读出表面积与
孔容；判读规则三者通用——过原点直线给出总表面积，下弯表示微孔填充，上弯表示毛细凝聚，
高段直线的斜率与截距给出外表面积与中孔（或微孔）孔容。

\textbf{$\alpha_s$ 方法}不用吸附层厚度，改用归一的吸附量 $a_s$，以非孔性参考吸附剂上
相对压力为 $0.4$ 时的吸附量 $n_{0.4}$ 作分母：

\begin{equation}\label{eq:as}
a_s=\frac{n\!\left(\ppzero\right)}{n_{0.4}}
\end{equation}

对同一类吸附剂建立统一的 $\alpha_s$ 曲线（$\alpha_s$ 对 $\ppzero$ 的标准数据表），再把
多孔样品各相对压力下的 $n$ 对 $a_s$ 作图得 $V$--$\alpha_s$ 图，判读方式与 $V$--$t$ 图
完全平行。\tocheck{未对照原页面图像：原书式(1-39) 的完整形式（分子分母记号，以及横轴是
$\ppzero$ 还是归一化后的 $(p/p_0)/0.4$）OCR 失真。}

\textbf{$n$ 方法}进一步认为吸附层厚度（以层数 $n=V/V_{\mathrm{m}}$ 表示）不仅取决于
相对压力，还取决于吸附质与吸附剂相互作用的强度；由于 BET 公式中的 $C$ 值恰好反映这一
强度，故按 $C$ 值分级（例如 $C>300$、$300>C>100$、$100>C>40$、$40>C>30$ 等）分别建立
$n$ 标准等温线。多孔样品作图时须先按自身 $C$ 值\textbf{选择合适的那一条}曲线；按 $C$ 值
分类的做法，实质是把 $t$ 方法「存在一条与吸附剂无关的标准等温线」的假设放宽为「同一
相互作用强度等级内存在一条标准等温线」\src{S1 §1.2.2 3)b，书内 p.53--54 / PDF p.72--73}。\tocheck{未对照原页面图像：原书
式(1-40) 的写法与 $C$ 值分级的完整档位（OCR 只保留 4 档，末档被截断）未能确认。}

\subsubsection{DFT/NLDFT：把分布写成积分方程}

前述方法都建立在若干假设之上，并有各自的适用孔径范围。基于分子统计力学的密度泛函理论
（DFT）与蒙特卡罗（MC）模拟不预设吸附机制，在分子水平描述吸附质相的构型，因而可以更
准确地计算吸附等温线。其中\textbf{非局域密度泛函理论}（NLDFT）算法相对简单、运算时间短、
精度较高，且\textbf{适用于微孔与中孔全范围}，是目前主流的孔尺寸分布算法。DFT 的基本
思想是：非均匀流体的自由能可表达为局部密度的泛函。NLDFT 采用巨正则系综，构筑限定空间内
流体分子的巨势 $\Omega$ 与局部密度分布 $\rho(r)$ 之间的函数关系：

\begin{equation}\label{eq:nldft-pot}
\Omega\!\left[\rho(r)\right]=F\!\left[\rho(r)\right]
+\int\!\rho(r)\left[V_{\mathrm{ext}}(r)-\mu\right]\mathrm{d}r
\end{equation}

式中 $F[\rho(r)]$ 为无外场时流体分子的亥姆霍兹自由能泛函，含理想气体自由能项与过量自由
能（excess free energy）项；过量自由能又分相吸与相斥两部分，相吸部分按 van der Waals 型
平均场近似（MFT）处理，相斥部分按非均匀硬球流体处理、多采用 Tarazona 的平滑密度近似；
$V_{\mathrm{ext}}(r)$ 为流体分子距孔壁 $r$ 处受到的壁面作用势，与所选孔形（裂隙、圆柱形
或球形）和孔宽有关；$\mu$ 为流体分子的化学势。热力学平衡时巨势对局部密度的导数为零，
迭代求解即得平衡密度分布 $\rho_{\mathrm{eq}}(r)$，进而算出吸附热与吸附等温线。由实验
等温线求分布时，先用统计力学方法计算该吸附质／吸附剂体系在\textbf{各单一孔宽 $D$} 上的
模型等温线 $N(\ppzero,D)$ 构成一套 \textbf{kernel}，再把实验等温线写成各孔贡献的积分

\begin{equation}\label{eq:gae}
N\!\left(\ppzero\right)=\int_{D_{\min}}^{D_{\max}}
N\!\left(\ppzero,D\right)f(D)\,\mathrm{d}D
\end{equation}

数值求解这一\textbf{积分吸附方程}（integral adsorption equation），即得孔尺寸分布函数
$f(D)$。\tocheck{未对照原页面图像：原书式(1-45)--(1-47) 的记号（巨势泛函的确切写法、
$V_{\mathrm{ext}}$ 的自变量、kernel 与积分限符号）在 OCR 中残缺，此处按 NLDFT 标准表述
重写，未逐符号核对。}

\begin{keycond}
\textbf{NLDFT 的成立条件与已知偏差。}（1）\textbf{kernel 必须与实验体系一致}：其数值
取决于孔形模型、吸附质分子之间（fluid--fluid）及吸附质分子与孔壁之间（fluid--solid）的
相互作用参数、以及权重函数等模型假设；用错 kernel，解出的分布会有显著错误。
（2）\textbf{低相对压力段常有系统性偏差}，分析中微孔复合材料时尤为明显；把相互作用参数
当可调参数处理可降低拟合偏差，但其取值必须能合理再现气液平衡密度与压力、气液表面张力等
流体性质，通常通过先拟合一个与样品表面化学结构相似的无孔固体上的标准等温线来优化。
（3）\textbf{孔壁在分子尺度光滑}的假设与实际不符；淬火固体密度泛函理论（QSDFT）引入
固体密度的不均匀性与各向异性，改善了多种微孔—中孔材料的表征精度。（4）\textbf{拟合好
不等于模型对}：式~\eqref{eq:gae} 拟合程度越高结果越可靠，但理论模型本身是对真实的理想化
概括，引入的假设与近似越多、物理意义越模糊，最终仍会退化为经验近似。规整孔道材料
（MCM-41、SBA-15）为校验提供了理想模型吸附剂，其计算结果已通过与 XRD、TEM 等直接观测
方法的对比得到验证\src{S1 §1.2.2 4)，书内 p.56--58 / PDF p.75--77}。
\end{keycond}

气体吸附法已标准化：GB/T 21650.2—2008《压汞法和气体吸附法测定固体材料孔尺寸分布和孔隙度
第 2 部分：气体吸附法分析介孔和大孔》，GB/T 21650.3—2011 同系列第 3 部分《气体吸附法
分析微孔》。

\subsubsection{压汞法}

\begin{principle}
汞对一般固体不润湿，欲使汞进入孔内必须施加外压。取刚性圆柱形孔模型（孔半径 $r$），抵抗
汞进入的界面张力沿孔壁圆周作用，量值为 $2\pi r\gamma\cos\theta$；克服它的外力（压强
$p$）作用在孔截面上，量值为 $\pi r^{2}p$。平衡时两力相等，得 \textbf{Washburn 方程}：

\begin{equation}\label{eq:washburn}
p=-\frac{2\gamma\cos\theta}{r}
\end{equation}

式中 $\gamma$ 为汞的表面张力，$\theta$ 为汞与固体表面的接触角。各类物质的接触角为
$135^{\circ}$--$150^{\circ}$，通常取平均值 $140^{\circ}$，$\cos140^{\circ}<0$，故
$p>0$，且\textbf{压入的孔半径与所受外压力成反比}。代入常用取值写成便于使用的形式

\begin{equation}\label{eq:washburn-num}
r/\unit{nm}=\frac{750}{p/\unit{MPa}}
\end{equation}

量纲检查：$2\gamma|\cos\theta|$ 为 $\unit{N.m^{-1}}$，除以 $p$（$\unit{Pa}=\unit{N.m^{-2}}$）
得 $\unit{m}$，与 $r$ 一致。汞填充孔的顺序是\textbf{先外部后内部、先大孔后中孔再小孔}；
测不同外压下进入孔内的汞量，即得相应孔宽的孔容。
\tocheck{未对照原页面图像：原书式(1-48)、(1-49) 的符号与常数取值（$\gamma$ 取 0.480 还是
0.485 N/m、$r$ 与 $p$ 的单位约定）OCR 失真；按 $\gamma=\qty{0.480}{N.m^{-1}}$、
$\theta=140^{\circ}$ 算得常数约 735，常用值 750 对应 $\gamma=\qty{0.485}{N.m^{-1}}$。}
\end{principle}

可测孔宽上下限由最低与最高填充压力决定：目前压汞仪使用压强最大 $200$--$400\,\unit{MPa}$，
可测孔半径范围约 $3\,\unit{nm}$ 至 $100$\,µm；由式~\eqref{eq:washburn-num} 可
自行校核——$p=\qty{200}{MPa}$ 对应 $r\approx\qty{3.7}{nm}$，与下限一致。压汞法原则上可测
直径大于 $3\,\unit{nm}$ 的开孔系统，但\textbf{一般只用于大孔}，中孔与微孔的分析用气体
吸附法。数据处理与气体吸附法同构：把进汞曲线（$V$--$p$）的压力坐标换算为孔尺寸即得积分
进汞孔容分布，求导得微分分布\src{S1 §1.2.2 2)3)，书内 p.58--60 / PDF p.77--79}。

\begin{measure}
\textbf{压汞法的实验条件。}样品池先抽空，再注满汞，过剩的汞返回储汞器；压力经液压油传给
汞，升压方式有连续扫描与步进扫描两种。随压力升高汞被压入孔内、汞液面下降，用浸入汞中的
电极检测汞体积随压力的变化，记录\textbf{进汞曲线}；压力升到预定值后仪器自动降压，记录
\textbf{退汞曲线}。压汞仪一般分低压系统（注汞、测大孔结构）与高压系统（测中孔大小与
分布）。为缩短分析时间并保护高压部件，\textbf{应按样品实际需要设定最高压力，不必都设到
仪器上限}。汞必须经酸洗、干燥和蒸馏，因为纯净度同时影响接触角与表面张力；还应做
\textbf{不装样品的空载实验}，以修正汞自身被压缩而贡献的虚假侵入体积，以及样品、样品管
和其他仪器元件带来的误差\src{S1 §1.2.2 2)，书内 p.59--60 / PDF p.78--79}。
\end{measure}

\subsubsection{回滞环与孔形判读}

吸附支与脱附支不重合，根源是发生毛细凝聚的中孔的几何结构：凝聚与蒸发时凹液面的情况
不同。因此\textbf{用吸附支与用脱附支分别计算孔尺寸分布，可能给出不同结果}；而且无论用哪
一支都需预先设定孔形状模型，算出的都是流体渗透等效孔的分布。催化研究关心的是能合理反映
孔结构对内扩散影响的结果，并不特别在意「真实」孔几何。

该书建议一般用\textbf{脱附支}数据计算，理由有五条：理想的两端开口圆柱形孔两支重合；此类
孔的吸附支对应圆柱形凹液面，脱附支对应孔口处的半球形凹液面，与 Kelvin 方程的标准形式
匹配；平板孔与狭缝孔在吸附时不发生毛细凝聚，只有脱附支反映真实孔隙；「墨水瓶」孔的吸附支
给出空腔尺寸、脱附支反映喉管尺寸，而内扩散速率恰受最窄喉管限制；吸附时毛细凝聚前可能需
一定过饱和，Kelvin 方程的热力学平衡未必达到，而脱附时孔内凝聚液与液体本体性质接近
\src{S1 §1.2.2 1)，书内 p.47 / PDF p.66}。

\begin{keycond}
\textbf{「用脱附支」不是无条件的。}$77\,\unit{K}$ 氮吸附下，脱附支能真实反映孔尺寸的
\textbf{下限是 $5$--$6\,\unit{nm}$}。对 H2(a) 型回滞环的「墨水瓶」孔，当喉管宽度小于这一
临界尺寸时，脱附支的毛细蒸发作用急剧下降、回滞环\textbf{强制闭合}，此时用脱附支算出的
只是这个临界尺寸，而不是真实分布。H3、H4、H5 型回滞环在 $\ppzero=0.4$--$0.5$ 区间也发生
强制闭合，用脱附支计算会在 $4$--$6\,\unit{nm}$ 之间出现一个\textbf{假峰}
\src{S1 §1.2.2 1)，书内 p.47--48 / PDF p.66--67}。
\end{keycond}

要了解中孔的连通情况，可做\textbf{回滞环扫描}（hysteresis scanning）：对比吸附--脱附扫描
曲线与边界吸附／脱附等温线的回滞环形状，从中提取孔连通性信息——这是 BJH 类「孔互不连通」
假设无法覆盖的部分。

\begin{figure}[htbp]
\centering
\begin{tikzpicture}
\begin{axis}[
  width=\onecolwidth, height=52mm,
  xmode=log, log basis x=10,
  xmin=0.8, xmax=60, ymin=0, ymax=1.02,
  xtick={1,2,5,10,20,50}, xticklabels={1,2,5,10,20,50},
  ytick={0,0.2,0.4,0.6,0.8,1.0},
  xlabel={孔宽 $D_p$ / nm},
  ylabel={$\mathrm{d}V_p/\mathrm{d}\lg D_p$ / $\unit{cm^3.g^{-1}}$},
  tick label style={font=\fontsize{8}{9.5}\selectfont},
  label style={font=\fontsize{8.5}{10}\selectfont},
  legend style={font=\fontsize{8}{9.5}\selectfont, at={(0.98,0.97)},
                anchor=north east, draw=none, fill=none, row sep=0.5pt},
  legend cell align=left]
% 峰形用有界函数 1/(1+((lg D-a)/w)^2)，避免小宽度指数导致 Dimension too large
\addplot[domain=0.8:60,samples=300,mainline,black]
  {0.62/(1+((log10(x)-log10(1.30))/0.090)^2)+0.34/(1+((log10(x)-log10(3.90))/0.140)^2)};
\addplot[domain=0.8:60,samples=300,thinline,black,dashed]
  {0.70/(1+((log10(x)-log10(4.70))/0.160)^2)+0.18/(1+((log10(x)-log10(1.95))/0.075)^2)};
\addplot[domain=0.8:60,samples=300,thinline,black,densely dotted]
  {0.55/(1+((log10(x)-log10(3.15))/0.130)^2)};
\legend{NLDFT（微孔＋中孔双峰）,BJH，脱附支,BJH，吸附支}
\draw[accentsoft,dashed,thinline] (axis cs:2,0) -- (axis cs:2,1.02);
\node[anchor=south west,font=\fontsize{8}{9.6}\selectfont,color=accent,
      text width=22mm,align=left] at (axis cs:2.15,0.02) {微孔上限 \qty{2}{nm}};
\end{axis}
\end{tikzpicture}
\caption{同一组等温线的两种反演结果。读者应看三处：\textbf{峰的位置}（BJH 用吸附支时中孔
峰在约 $\qty{3.2}{nm}$，用脱附支时右移到约 $\qty{4.7}{nm}$，而 NLDFT 给出约
$\qty{3.9}{nm}$）；\textbf{峰的个数}（只有 NLDFT 解出约 $\qty{1.3}{nm}$ 的微孔峰，BJH
在微孔区给不出真实信息，却可能在约 $\qty{1.9}{nm}$ 处留下一个由外推造成的虚假小峰）；
\textbf{峰高}（中孔孔容的估计值随模型变化）。三条曲线出自同一组数据，差别完全来自模型，
因此任何孔尺寸分布都必须连同模型一起报告。本图按 BJH 与 NLDFT 的典型反演行为构造，是
说明模型选择影响的示意曲线，不代表任何具体样品的实验数据。\modelmark
\figsrc{据 S1（上册）§1.2.2，书内 p.45--48 / PDF p.64--67 所述模型差异计算绘制。}}
\label{fig:psd-model}
\end{figure}

\begin{table}[htbp]
\centering
\caption{各类孔的适用分析方法与不适用方法}
\label{tab:pore-class}
{\srcsize
\begin{tabularx}{\columnwidth}{@{}llXX@{}}
\toprule
孔宽 & 名称 & 适用模型／方法 & 不适用或必须慎用（原因） \\
\midrule
$<\qty{2}{nm}$ & 微孔 &
D-R／D-A 方程、HK／SF 法、MP 法（$V$--$t$）、NLDFT &
BJH 法：Kelvin 方程在微孔内不成立，会给出 $0.5$--$2\,\unit{nm}$ 的虚假峰；
BET 面积只能称等效表面积 \\
$\qtyrange{2}{50}{nm}$ & 中孔 &
BJH／BdB、$t$ 图法、$\alpha_s$ 法、$n$ 法、NLDFT &
压汞法：只在低压段少量覆盖，小孔区受高压与接触角假设支配；
HK／SF：势能模型只对微孔标定 \\
$>\qty{50}{nm}$ & 大孔 &
压汞法、电子显微镜等直接观测法 &
氮吸附法：$\ppzero$ 接近 $1$，凝聚量测不准，Kelvin 方程不可靠 \\
\bottomrule
\end{tabularx}}
\end{table}

\begin{measure}
\textbf{气体吸附法测孔分布的条件。}（1）\textbf{吸附质选择}：$\qty{77}{K}$ 氮吸附是测
BET 表面积和中孔孔尺寸分布最常用的体系，但\textbf{不一定是测微孔孔容的最好选择}——氮
分子与分子筛孔道内阳离子之间存在较强的四极矩诱导相互作用，影响分析结果的真实性；IUPAC
建议改用惰性的球形分子 Ar（$87\,\unit{K}$）分析微孔。（2）\textbf{超低压段必须测准}：
常规吸附测量从 $\ppzero\approx0.01$ 开始，此时大部分微孔可能已经充满，因而测不到微孔
分布；高分辨低压吸附技术可从 $10^{-6}$ 甚至 $10^{-7}$ 的相对压力开始，这要求仪器有高
真空度和高精度压力传感器。（3）\textbf{脱气条件更苛刻}：排除微孔内的吸附物比中孔、大孔
困难得多。（4）\textbf{氦气截留}（helium entrapment）：商品自动吸附仪在液氮温度下用氦气
测死体积，微孔样品在此过程中会吸附相当量的氦且难以抽净，部分氦在随后吸附氮时释放，使超
低压段平衡压力偏高。解决办法是把死体积测量与微孔吸附测定分成两个独立过程：死体积测完后
重新加热脱氦，再把已测死体积作为已知参数输入仪器。（5）\textbf{避免回充气}：预处理后的
气体回充同样会使超低压段平衡压力偏高，应额外脱气或选不回填气体的分析方式。
（6）\textbf{平衡时间长}：微孔样品的平衡时间明显长于中孔、大孔样品，需放宽平衡判据
\src{S1 §1.2.2 2)，书内 p.49 / PDF p.68}。
\end{measure}

\begin{dataflow}
从一条等温线到一份可复核的孔尺寸分布：
（1）\textbf{定孔容}。纯微孔吸附剂的孔容对应 I 型等温线的平台吸附量；中微孔混合吸附剂的
微孔填充在 $\ppzero=0.01$--$0.15$ 之间完成，但此时总吸附量还含外表面上的亚单层吸附量；
IV 型等温线表明含中孔，高 $\ppzero$ 段的终止平台吸附量即中孔与更大孔的总孔容，若平台
不易分辨，总孔容通常取 $\ppzero=0.95$ 处的吸附量，对应孔宽上限 $\qty{50}{nm}$。孔容
除以样品质量得比孔容。
（2）\textbf{选模型}。微孔段用 D-R／HK／SF／MP／NLDFT，中孔段用 BJH／BdB／$t$ 图／
NLDFT，大孔段用压汞法；先由等温线类型与回滞环类型缩小范围。
（3）\textbf{反演}。按式~\eqref{eq:halsey} 或标准曲线算 $t$，按式~\eqref{eq:kelvin} 算
$r_k$，再按式~\eqref{eq:bjh-balance} 递推或按式~\eqref{eq:gae} 解积分方程，得积分分布，
求导得微分分布，取峰位为最可几孔宽。
（4）\textbf{交叉验证}。同一组数据换支（吸附支与脱附支各算一遍）、换模型、必要时换吸附质
各算一遍；分布差异大的孔宽区间不得写进结论。
（5）\textbf{报告}：按本节开头的要求，连同实验方法、$\ppzero$ 窗口与模型（含 kernel）名称
一并给出。
\end{dataflow}

\begin{pitfall}
\textbf{（1）用 BJH 分析微孔材料会得到虚假孔径峰。}高硅铝比 MFI 分子筛的高分辨低压吸附
在微孔填充段常出现台阶状跃升：骨架与氮、氩的强相互作用使微孔内吸附质由无序流体相转变为
类晶体固相，每个 MFI 单胞容纳的吸附质分子数由 23 个增至 30 个；骨架硅铝比越高、表面组成
越均匀，现象越明显（如纯硅的 silicalite-1）。此时若用 SF 法处理氩的等温线，
\textbf{或不正确地用 BJH 法处理氮的等温线（BJH 法不适用于该 $\ppzero$ 区间）}，都会得到
孔宽集中在 $0.5$--$2\,\unit{nm}$ 的虚假孔尺寸分布峰\src{S1 §1.2.2 2)，书内 p.51 / PDF p.70}。

\textbf{（2）把中孔材料的强吸附误判为微孔填充。}比表面积非常大的规整中孔材料单层吸附量
极大，低相对压力下就有显著吸附量；若把这部分吸附当成微孔填充再用 HK 或 SF 法分析，会
\textbf{人为制造出一个 $\qty{1}{nm}$ 左右、相当宽的孔尺寸分布}，而同一组数据用 $t$ 图法
作图仍是一条过原点的直线，说明它是纯中孔材料。两条判据矛盾时，以过原点直线为准，并换用
Ar 吸附复测。

\textbf{（3）脱附支在 $\ppzero\approx0.42$ 附近的伪峰。}H2(a) 型（墨水瓶孔）的回滞环在
喉管宽度小于临界尺寸时强制闭合；H3、H4、H5 型回滞环在 $\ppzero=0.4$--$0.5$ 区间同样
强制闭合。用脱附支计算时，分布曲线在 $4$--$6\,\unit{nm}$ 之间出现一个\textbf{假峰}，其
物理来源是\textbf{拉伸强度效应}（tensile strength effect）——凝聚液在张力下失稳蒸发，
而不是孔内液体真的在对应曲率下蒸发。因此\textbf{当脱附支给出的分布集中在小于
$\qty{12}{nm}$、特别是 $4$--$6\,\unit{nm}$ 时一定要慎重}：应改用吸附支对比，或采用 NLDFT
分析，还可以改变吸附温度甚至更换吸附质，最好再用直接观测法（如 TEM）确认
\src{S1 §1.2.2 1)，书内 p.47--48 / PDF p.66--67}。

\textbf{（4）压汞法在小孔区由压力与接触角假设支配。}接触角取平均 $140^{\circ}$ 在许多
情况下过于粗糙，由此导致的孔尺寸分析误差可达 $50\%$，准确结果要求汞与每种样品的接触角
实测。高压段还有两个额外贡献：汞自身被压缩的体积，以及可压缩样品在
$200$--$400\,\unit{MPa}$ 下被压紧甚至破坏。此外真实孔道常带窄喉管，汞只有达到喉管入口
半径对应的压力才能充满整个孔，算出的分布偏小于实际空腔尺寸；退汞时喉管又造成汞截留与
滞后，使退汞曲线总在进汞曲线上方。第一次进汞的低压段进汞量主要补充颗粒堆积体的体积压缩
与颗粒间隙，退汞时这部分汞不被记录，因此\textbf{第一次进汞曲线与退汞曲线不闭合，需反复
几次才闭合}。与汞形成汞齐的金属样品不适用
\src{S1 §1.2.2 2)3)4)，书内 p.58--60 / PDF p.77--79}。

\textbf{（5）不同模型给出的平均孔径不可直接互比。}各方法的孔宽定义、孔形假设（圆柱形／
平板／球形）、流体性质与积分区间都不同，得到的都是\textbf{等效模型孔的名义值}。把 BJH
脱附支的平均孔径与 NLDFT 的最可几孔宽、压汞法的孔半径写在同一张表里做趋势比较，是把模型
差异当成了材料差异；可比的前提是同一方法、同一模型、同一吸附质。

\textbf{（6）把 $t$ 图截距当作总孔容。}$V$--$t$ 曲线高段的截距只给出\textbf{中孔孔容}
（微孔吸附剂则为微孔孔容），不含大孔；大孔孔容需由压汞法或密度法给出。
\end{pitfall}

\begin{crosscheck}
\textbf{（1）与直接观测法互查。}校验 $t$ 方法与 NLDFT 最有力的办法，是用规整孔道结构的
多孔吸附剂（MCM-41、SBA-15）把直接观测结果与间接分析结果对比；对不规整孔材料，用透射
电子显微镜抽查孔宽量级，即可发现分布的明显错误（如把中孔材料报成微孔材料）。
\textbf{（2）与压汞法互查。}两法在约 $3$--$50\,\unit{nm}$ 区间重叠：吸附法在中孔段更
精细，压汞法在大孔段不可替代。若两者在同一孔宽区间给出的孔容相差数倍，应先复查预处理
条件与接触角取值，而不是取平均。
\textbf{（3）与密度、孔隙率互查。}总孔容应与由骨架密度和表观密度算出的孔隙率自洽（见下一
节）；不自洽说明孔容的积分区间没有覆盖全孔。
\textbf{（4）与表面积互查。}由孔尺寸分布积分得到的内表面积应与 BET 等效表面积同量级；
若分布积分面积远大于 BET 面积，通常意味着模型中掺入了虚假的小孔峰。
\textbf{（5）与反应数据互查。}内扩散速率受最窄喉管尺寸限制，因此分布必须能解释扩散实验
（如有效因子、Thiele 模数）的结果；若分布显示孔很大而扩散限制很强，优先怀疑模型给出了
错误的孔口尺寸，而不是怀疑扩散实验
\src{S1 §1.2.2，书内 p.44--45、p.47 / PDF p.63--64、p.66}。
\end{crosscheck}

\srcnote{本节来源：S1（上册）§1.2.2，书内 p.43--60 / PDF p.62--79。
式~\eqref{eq:kelvin}--\eqref{eq:washburn-num} 依 OCR 正文给出的变量定义与学科标准形式
重建，未逐式对照原页面图像；各式的成立条件与量纲自检见相应 \texttt{keycond} 与正文。
图~\ref{fig:psd-model} 为按模型差异计算的教学示意曲线，表~\ref{tab:pore-class} 为据正文
归纳，均非原书图表。压力--孔宽对照表与 HK／SF 法所用的吸附剂、吸附质参数见原书表 1-7
至表 1-9，本分片未逐值转录。}
